% !TeX root = ../thuthesis-example.tex

\chapter{选题背景及意义}

\section{量子互联网的发展愿景与核心挑战}

量子互联网的核心目标，是将地理上分离的量子处理器、量子存储器与量子传感器通过量子链路互联，形成能够支撑量子隐形传态、分布式量子计算、量子密钥分发和量子传感等任务的网络基础设施。实现这些应用的基础资源是端到端纠缠对，而纠缠对的保真度直接决定上层量子任务的可靠性。因此，如何在大规模、多跳网络中高效分发高保真纠缠资源，是构建实用量子互联网必须解决的网络层问题。

与经典互联网不同，量子网络中的路径质量不能只由跳数、带宽或静态时延描述。纠缠生成具有概率性，链路保真度会受到信道噪声影响，量子态在存储器中等待时还会发生退相干。尤其在多跳纠缠交换中，较早建立完成的链路必须等待较慢链路完成；等待时间由路径内部的时间瓶颈决定，并会通过后续纠缠交换传播到端到端保真度。因此，量子路由必须把链路初始保真度、链路建立时间、节点相干时间和交换过程放到统一模型中考虑。

\section{量子路由的非加性困难}

经典最短路径算法通常依赖加性链路代价：一条路径的总代价可以由每条边的代价相加得到，局部更优的前缀在拼接相同后缀后仍然保持不劣。然而，本文当前实现采用的量子路由度量并不满足这一结构。一个候选路径 label 会保留整条前缀上所有链路段的物理参数；每当扩展一条新边时，系统重新计算整条前缀的 $T_{\max}$、各段等待时间、等待后的 Bell 对角态以及递归纠缠交换结果。

这意味着路径质量不是一个可加标量。一个在中间节点处保真度较高的前缀，可能因为携带较大的 $T_{\max}$，在接入同一快速后缀后使后缀链路等待更久，从而被退相干反超。该现象破坏了单标签 Dijkstra 的局部剪枝前提，也是本文重点讨论的“前缀陷阱”问题。

\section{qns-3平台与本文工作}

qns-3 是基于 ns-3 的量子-经典混合网络仿真平台，其特点是将量子操作、经典信令和网络协议事件纳入同一离散事件时间轴。本文工作在 qns-3 中补充并分析量子链路层与量子网络层，重点包括：

\begin{enumerate}
  \item 建立与当前代码一致的路径级保真度模型，在路由标签中显式保存链路段参数，并通过 $T_{\max}$ 对齐、相位翻转存储退相干和 Bell 对角态交换计算 \texttt{predicted\_fidelity}；
  \item 实现并比较单标签 \texttt{DijkstraRoutingProtocol} 与多标签 \texttt{SlicedDijkstraRoutingProtocol}，将前者作为 $K=1$ best-first 基线，将后者作为沿 $T_{\max}$ 轴保留多个候选前缀的切片算法；
  \item 证明当前路由度量不满足路径拼接保序性，说明单标签剪枝为什么可能丢掉未来更优的前缀，并讨论切片算法能够缓解但不能完全消除该问题的边界；
  \item 为后续真实拓扑和混态仿真实验保留评估框架，后续实验将重点检查切片宽度、标签上限和最终端到端保真度之间的关系。
\end{enumerate}

因此，本文的重点不是把已有实验结果包装成结论，而是先把物理建模、算法语义和理论证明统一到当前实现，为后续实验设计提供可靠基础。
